\begin{table}[H]
    \centering
    \caption{AUC no conjunto de teste sem e com as variáveis do município de residência.}
    \label{tab-res-cenarios}
    \footnotesize
    \renewcommand{\arraystretch}{1.2}
    \setlength{\tabcolsep}{4pt}
    \begin{tabular}{l c c c c}
        \hline
        \textbf{Modelo} & \textbf{Sem municipais} & \textbf{Com municipais} & \textbf{Diferença (IC 95\%)} & \textbf{p} \\
        \hline
        Regressão Logística & 0,698 (0,691–0,705) & 0,707 (0,700–0,713) & 0,008 (0,006–0,011) & < 0,001 \\
        Árvore de Decisão & 0,696 (0,689–0,703) & 0,701 (0,694–0,707) & 0,004 (0,001–0,008) & 0,022 \\
        Random Forest & 0,702 (0,695–0,709) & 0,716 (0,709–0,723) & 0,014 (0,011–0,017) & < 0,001 \\
        MLP & 0,701 (0,694–0,708) & 0,713 (0,706–0,720) & 0,012 (0,009–0,015) & < 0,001 \\
        LightGBM & 0,701 (0,694–0,708) & 0,724 (0,718–0,731) & 0,023 (0,019–0,027) & < 0,001 \\
        \hline
    \end{tabular}
    \legend{Nota: AUC (IC 95\%). Sem municipais = apenas variáveis individuais; com municipais = variáveis individuais e do município de residência; diferença = com municipais menos sem municipais. IC e p pelo teste de DeLong para curvas ROC pareadas.\\Fonte: elaborado pelo autor.}
\end{table}
